\begin{frame}{RNN의 근본적 한계: 순차성}
  \begin{itemize}
    \item 게이트를 추가해도 RNN은 여전히 \textbf{순차적으로 한
      시점씩} 처리해야 한다 --- 100번째 단어를 처리하려면
      1번째부터 99번째까지 \textbf{순서대로 다 거쳐}야 함
    \item 이 \textbf{순차성} 때문에:
    \begin{itemize}
      \item \textbf{병렬화가 어려}움
      \item 아주 긴 시퀀스에서는 \textbf{여전히} 먼 과거 정보가
        흐려진다
    \end{itemize}
  \end{itemize}
  \vskip0.6em
  \begin{block}{다음 장으로}
    Week 12의 \kb{Attention/Transformer}는 이 \textbf{순차성 자체를
    없애는} 완전히 다른 접근이다 --- ``필요할 때마다 과거 전체를
    직접 다시 들여다본다''.
  \end{block}
  \vskip0.5em
  \kq{RNN은 ``순서가 있는 데이터를 다루려면 과거를 기억하는
  상태가 필요하다''는 통찰의 \textbf{가장 단순한 구현}이다 ---
  다음 장은 이 구조의 한계를 극복하는 이야기다.}
\end{frame}
